The proposed framework is validated on two representative cases with contrasting energy demand characteristics: a heating-dominated German community and a cooling-dominated campus in Songshan Lake, China. Four groups of experiments are designed to systematically evaluate the two contributions of this work.

\subsection{Case descriptions}
\label{sec:case_desc}

\subsubsection{Case 1: German community (heating-dominated)}

The first case represents a distributed energy community in Germany with significant heating demand and moderate cooling requirements. The hourly load and weather data span a full year (8760 hours) and include electricity demand, heating demand, cooling demand, solar irradiance, wind speed, and ambient temperature. The electricity tariff is \SI{0.0025}{\text{\euro}/kWh} (flat rate), the natural gas price is \SI{0.0286}{\text{\euro}/kWh}, and the capacity charge is \SI{114.29}{\text{\euro}/(kW \cdot a)}. Both wind and solar resources are available at this site.

\subsubsection{Case 2: Songshan Lake campus (cooling-dominated)}

The second case represents a university campus in Songshan Lake, Dongguan, China, located in a subtropical climate zone. The annual cooling load (\SI{2.91}{GWh}) significantly exceeds the heating load, with a peak cooling demand of \SI{5797}{kW}. The electricity price is \SI{0.6746}{\text{\textyen}/kWh} (uniform research rate), the gas price is \SI{0.45}{\text{\textyen}/kWh}, and no capacity charge applies. Wind resources are negligible at this site ($P_{\text{wt}}^{\text{ub}} = 0$), so only solar generation is considered.

\Cref{tab:case_params} summarizes the key parameters for both cases.

\begin{table}[htbp]
\centering
\caption{Key parameters for the two case studies.}
\label{tab:case_params}
\small
\begin{tabular}{@{}llcc@{}}
\toprule
Parameter & Unit & German & Songshan Lake \\
\midrule
\multicolumn{4}{@{}l}{\textit{Energy prices}} \\
Electricity price & \euro/kWh or \textyen/kWh & 0.0025 & 0.6746 \\
Gas price & \euro/kWh or \textyen/kWh & 0.0286 & 0.45 \\
Capacity charge & \euro/(kW$\cdot$a) or \textyen/(kW$\cdot$a) & 114.29 & 0 \\
\midrule
\multicolumn{4}{@{}l}{\textit{Equipment upper bounds (kW)}} \\
PV & kW & 10{,}000 & 1{,}000 \\
WT & kW & 10{,}000 & 0 \\
GT & kW & 10{,}000 & 800 \\
EHP & kW & 3{,}000 & 300 \\
EC & kW & 1{,}000 & 5{,}000 \\
AC & kW & 1{,}000 & 1{,}500 \\
ES & kW & 20{,}000 & 2{,}000 \\
HS & kW & 6{,}000 & 500 \\
CS & kW & 2{,}000 & 3{,}000 \\
\midrule
\multicolumn{4}{@{}l}{\textit{Carnot battery (when enabled)}} \\
CB power & kW & 2{,}000 & 500 \\
CB capacity & kWh & 10{,}000 & 3{,}000 \\
CB round-trip efficiency & -- & 0.60 & 0.60 \\
\midrule
\multicolumn{4}{@{}l}{\textit{Equipment efficiencies}} \\
GT electrical efficiency & -- & 0.30 & 0.30 \\
GT thermal efficiency & -- & 0.45 & 0.45 \\
EHP COP & -- & 3.5 & 3.5 \\
EC COP & -- & 4.0 & 4.0 \\
AC COP & -- & 0.7 & 0.7 \\
\midrule
\multicolumn{4}{@{}l}{\textit{Energy quality coefficients (Carnot exergy factors)}} \\
$T_0$ (ambient) & K & 298.15 & 300.15 \\
$T_h$ (heating supply) & K & 343.15 & 333.15 \\
$T_c$ (cooling supply) & K & 280.15 & 280.15 \\
$\lambda_e$ (electricity) & -- & 1.000 & 1.000 \\
$\lambda_h$ (heating) & -- & 0.131 & 0.099 \\
$\lambda_c$ (cooling) & -- & 0.064 & 0.071 \\
\bottomrule
\end{tabular}
\end{table}

\subsection{Experimental design}
\label{sec:exp_design}

Four groups of experiments are designed to validate the two contributions of this work, as summarized in \Cref{tab:experiments}.

\begin{table}[htbp]
\centering
\caption{Experimental design overview.}
\label{tab:experiments}
\small
\begin{tabular}{@{}cllp{5.5cm}@{}}
\toprule
Exp. & Case & Methods & Purpose \\
\midrule
1 & German & A/B/C/D/E & Multi-method comparison (Contribution 1) \\
2 & Songshan Lake & B vs.\ C (no CB) & Cross-climate validation \\
3 & Songshan Lake & C (no CB) vs.\ C (with CB) & Carnot battery ablation (Contribution 2) \\
4 & Songshan Lake & B vs.\ C (both with CB) & Joint validation of both contributions \\
\bottomrule
\end{tabular}
\end{table}

The five methods compared in Experiment~1 are:
\begin{itemize}
    \item \textbf{Method A} (Economic): Single-objective cost minimization, serving as the baseline.
    \item \textbf{Method B} (Std): Bi-objective with the net load volatility index (\Cref{eq:mi_std}), representing the conventional approach.
    \item \textbf{Method C} (EQD): Bi-objective with the proposed energy-quality-weighted Euclidean distance (\Cref{eq:mi_eqd}).
    \item \textbf{Method D} (Pearson): Bi-objective with the Pearson correlation index (\Cref{eq:mi_pearson}).
    \item \textbf{Method E} (SSR): Bi-objective with the self-sufficiency rate index (\Cref{eq:mi_ssr}).
\end{itemize}

All bi-objective experiments use NSGA-II with a population size of $N_{\text{ind}} = 50$ and $G_{\text{max}} = 100$ generations. The single-objective Method~A uses a DE-based genetic algorithm with the same population size and generation count. The 14 typical days are pre-computed by $k$-medoids clustering on the combined load and weather profiles.
